\noindent
\textbf{Conventions.}
 We work throughout over $\SC$. In particular, all schemes are $\SC$-schemes and all tensor products are taken over $\SC$ unless otherwise indicated. We often use the following partial order on dimension vectors: for any $n\in \mathbb{N}$ and for $u=(u_i),v=(v_i)\in \mathbb{Z}^n$, we define $u\le v$ if $u_i\le v_i$ for all $1\leq i\leq n$. We adopt the sign convention of King~\cite{king1994moduli} for $\theta$-stability, so a~module $M$ over an algebra is $\theta$-stable (resp.\ $\theta$-semistable) if $\theta(M)=0$ and every nonzero proper submodule $N\subsetneq M$ satisfies $\theta(N)>0$ (resp.\ $\theta(N)\ge 0$).

\section{Preliminaries}\label{sec:prelim}

\subsection{Three algebras arising from the McKay graph}\label{subsec:mckay}
Let $\Gamma\subset \GL(V)$ be a finite subgroup of the general linear group for an $m$-dimensional $\SC$-vector space $V$.
Choosing a basis of $V$, we can write the symmetric algebra of the dual vector space $V^*$ as a polynomial ring $R=\SC[x_1,\dots,x_m]$. The group $\Gamma$ acts dually on $R$: for $g\in \Gamma$ and $f\in R$, we have $(g\cdot f)(v) = f\big(g^{-1} v\big)$ for all $v\in V$. Let $R^\Gamma$ denote the $\Gamma$-invariant subring of~$R$.

List the irreducible representations of $\Gamma$ as $\{\rho_0, \rho_1, \dots, \rho_r\}$, where $\rho_0$ is the trivial representation.
The decomposition of $R$ as an $R^\Gamma$-module can be written \cite[Proposition~4.1.15]{goodman2009symmetry}
 \begin{gather*} %\label{eqn:decompR}
 R\cong \bigoplus_{0\leq i\leq r} R_i\otimes_\mathbb{C} \rho_i,
 \end{gather*}
 where $R_i\coloneqq \Hom_\Gamma(\rho_i,R)$ is an
 $R^\Gamma$-module; note that $R_0=R^{\Gamma}$.

From now on, let $\Gamma \subset \SL(2,\SC)$ be a finite subgroup acting as above on $R=\SC[x,y]$. The \emph{McKay graph} of $\Gamma$
has vertex set $\{0,1, \dots, r\}$ corresponding to the irreducible representations of~$\Gamma$, where for each $0\leq i, j\leq r$, we have $\dim\Hom_{\Gamma}(\rho_j , \rho_i \otimes V )$ edges joining vertices~$i$ and~$j$. McKay~\cite{McKay80} observed that this graph is an affine Dynkin diagram of ADE type. We frame this graph by introducing an additional vertex, denoted~$\infty$, together with an edge joining the vertices~$\infty$ and~$0$; we refer to the resulting graph as the \emph{framed McKay graph} of~$\Gamma$.

We now associate a doubled quiver to each graph. For this, consider
the set of pairs $Q_1$ comprising an edge in the framed McKay graph
and an orientation of the edge. For each $a\in Q_1$, we write~$t(a)$,
$h(a)$ for the vertices at the tail and head respectively of
the oriented edge, and we write $a^{\ast}$ for the same edge with
the opposite orientation. Define the \emph{framed McKay quiver of $\Gamma$},
denoted $Q$, to be the quiver with vertex set $Q_0= \{\infty,0,\dots ,r\}$
and arrow set $Q_1$. The (unframed) \emph{McKay quiver},
denoted $Q_\Gamma$, is the complete subquiver of $Q$ on
the vertex set $\{0, 1,\dots, r\}$. Using these quivers
we define the following three algebras.

First, let $\SC Q$ denote the path algebra of the framed McKay quiver $Q$. For $i \in Q_0$, let $e_i \in \SC Q$ denote the idempotent corresponding to the trivial path at vertex $i$. Let $\epsilon\colon Q_1 \to \{\pm 1\}$ be any map such that $\epsilon(a) \neq \epsilon(a^{\ast})$ for all $a \in Q_1$. The \emph{preprojective algebra of $Q$}, denoted $\Pi$, is defined as the quotient of $\SC Q$ by the ideal
\begin{equation}
\label{eqn:preprojectiverelation}
\left\langle \sum_{\head(a)=i} \epsilon(a) aa^{\ast} \mid {i \in Q_0}
\right\rangle.
\end{equation}
 The preprojective algebra $\Pi$ does not depend on the choice of the map $\epsilon$. For each $i\in Q_0$, we simply write $e_i\in \Pi$ for the image of the corresponding vertex idempotent.

Second, in the framed McKay quiver $Q$, let $b^*\in Q_1$ be the unique arrow with head at vertex~$\infty$. Define a $\SC$-algebra as the quotient of the preprojective algebra $\Pi$ by the two-sided ideal generated by the class of $b^*$:
 \[
 \PiQ\coloneqq \Pi/(b^*).
 \]
 Equivalently, if we define a quiver $Q^*$ to have vertex set $Q^*_0=\{\infty,0,1,\dots,r\}$ and arrow set $Q_1^*=Q_1\setminus \{b^*\}$, then $\PiQ$ is the quotient of the path algebra $\SC Q^*$ by the ideal of relations
 \begin{equation*}% \label{eqn:Arelations}
 \left\langle \sum_{\head(a) = i} \epsilon(a) aa^* \mid 0\leq i,
 \head(a^*)\leq r \right\rangle.
 \end{equation*}

 Third, let $\SC Q_\Gamma$ denote the path algebra of the (unframed) McKay quiver $Q_\Gamma$. Since $\SC Q_\Gamma$ is a subalgebra of $\SC Q$, we use the same symbol $e_i\in \SC Q_\Gamma$ for the idempotent corresponding to the trivial path at vertex $i$ of $Q_\Gamma$. The \emph{preprojective algebra of $Q_\Gamma$}, denoted $\PiGamma$, is defined as the quotient of $\SC Q_\Gamma$ by an ideal defined similarly to that from~\eqref{eqn:preprojectiverelation} for the quiver $Q_\Gamma$.

\begin{Remark}
 For each vertex $0\leq i\leq r$ we denote by $e_i$ the class in $\PiQ$ of the idempotent at the vertex $i$. Note that $\PiGamma$ is not a subalgebra of $\Pi$. On the other hand, $\PiGamma$ is isomorphic to the subalgebra $\oplus_{0 \leq i,j \leq r} e_j\PiQ e_i$ of $\PiQ$ due to \cite[Lemma~3.2]{craw2019punctual}, and moreover
 \begin{equation} \label{eq:PiQdecomp}
 \PiQ \cong \PiGamma \oplus \PiGamma b \oplus \SC e_\infty.
\end{equation}
as complex vector spaces. In particular, the classes $e_i$, $0 \leq i \leq r$ are also contained in $\PiGamma$.
\end{Remark}
